\documentclass[12pt]{letter}

\usepackage[T1]{fontenc}
\usepackage[utf8]{inputenc}
\usepackage{newtxtext}
\usepackage{newtxmath}
\usepackage[margin=1in]{geometry}

\signature{Yash Bisht\\Department of Mechanical and Industrial Engineering\\Indian Institute of Technology Roorkee}
\address{Department of Mechanical and Industrial Engineering\\Indian Institute of Technology Roorkee\\Roorkee 247667, Uttarakhand, India}

\begin{document}

\begin{letter}{The Editor-in-Chief\\Journal of Environmental Chemical Engineering}

\opening{Dear Editor,}

Please consider our manuscript entitled ``Source-aware machine learning for heavy-metal adsorption: distinguishing high-accuracy interpolation from cross-source generalization'' for publication in the \textit{Journal of Environmental Chemical Engineering}.

The manuscript addresses an important methodological issue in adsorption machine learning: high reported removal-efficiency prediction accuracy can depend strongly on whether the model is evaluated within familiar experimental sources or tested on unseen references, adsorbents, and source domains. Using a cleaned heavy-metal adsorption dataset of 563 records, we identify the neutral/low-concentration regime as a sparse applicability region with only 26 observations and perform primary modeling on the remaining 537 records. Source-specific XGBoost models reach $R^2$ up to 0.992, while restricted-domain reference-held-out screening reaches $R^2 = 0.305$. The work therefore provides a practical dual-track validation framework that separates high-accuracy interpolation, capacity-inclusive process monitoring, capacity-free pre-experiment screening, and external-source generalization.

We believe the manuscript fits the scope of the journal because it combines environmental treatment data, adsorption-process modeling, and reproducible chemical-engineering machine learning. The study is intended to improve how adsorption ML models are validated and reported, helping readers distinguish useful within-system prediction from unsupported claims of universal deployment.

The manuscript is original, has not been published previously, and is not under consideration elsewhere. The authors declare no known competing financial interests or personal relationships that could have appeared to influence the work reported in the manuscript.

\closing{Sincerely,}

\end{letter}

\end{document}
